% Ad Familiares 16.27 --- headnote
% ad-familiares-16-27, c. 15 Dec 44 BC, place uncertain

\textit{Quintus Cicero --- the orator's younger brother --- to
Tiro, written from an uncertain location at the end of December
44 BC, per the Perseus dateline \textit{loco incerto ex. m. Dec.
a.~710 (44)}. The salutation \textit{Q. CICERO TIRONI SVO S. P.
D.} identifies the sender as Quintus (not his son and namesake);
the warmth of the close (\textit{te fero in oculis}) marks the
long friendship between Quintus and his brother's freedman.}

\textit{The voice is unmistakably Quintus: blunter than his
brother, more openly contemptuous of the political leadership,
and given to scathing characterisations. The consuls-designate
for 43 are Aulus Hirtius and Vibius Pansa, dismissed here as
unmanned voluptuaries unfit for power. The ``bandit'' is Antony.
The Caesena reference and the Cossutian shops are obscure jabs
--- proverbial bywords for utterly worthless property, the point
being that the consuls-designate are not even fit to be put in
charge of \textit{those}. The closing image of kissing Tiro's
eyes if he meets him in the middle of the Forum is characteristic
unguarded Quintus.}
